\chapter{Mediastinoscopy}

\section*{What Is This Surgery?}
Mediastinoscopy is a minimally invasive surgical procedure performed to visually inspect and obtain tissue biopsies from the mediastinum, the central compartment of the chest located between the lungs. This surgery is primarily utilized for the diagnostic evaluation and staging of various mediastinal pathologies, most notably non-small cell lung cancer (NSCLC). By providing direct access to mediastinal lymph nodes, particularly those in the paratracheal and subcarinal regions, it allows for accurate histological assessment crucial for treatment planning. Beyond cancer staging, mediastinoscopy is also invaluable in diagnosing other conditions such as sarcoidosis, lymphoma, tuberculosis, and other mediastinal masses of unknown etiology, offering definitive tissue diagnosis where less invasive methods may be inconclusive. The procedure's ability to provide robust tissue samples makes it a cornerstone in the management of complex thoracic diseases, guiding oncological and medical therapies.

\section*{Alternative Names}
\begin{itemize}
  \item Cervical Mediastinoscopy
  \item Standard Mediastinoscopy
  \item Mediastinal Lymph Node Biopsy
  \item Carlens' Mediastinoscopy
  \item Pretracheal Mediastinoscopy
  \item Diagnostic Mediastinoscopy
\end{itemize}

\section*{Relevant Surgical Anatomy}
The mediastinum is the central compartment of the thoracic cavity, bounded by the sternum anteriorly, the vertebral column posteriorly, the pleura laterally, and the diaphragm inferiorly. It is conventionally divided into superior and inferior mediastinum, with the latter further subdivided into anterior, middle, and posterior compartments. Mediastinoscopy primarily accesses the superior and anterior aspects of the middle mediastinum. Key structures within the surgical field include the trachea, esophagus, and major vascular structures.

The trachea lies centrally, serving as the primary landmark. Anterior to the trachea, the innominate (brachiocephalic) artery crosses obliquely from left to right, often marking the inferior limit of safe dissection in the superior mediastinum. The superior vena cava (SVC) lies to the right of the trachea, while the aortic arch and its major branches (innominate artery, left common carotid artery, left subclavian artery) are to the left and anterior. The pulmonary artery bifurcation is located inferiorly, often near the subcarinal lymph node station. Important nerves include the recurrent laryngeal nerves, which ascend in the tracheoesophageal grooves (the left nerve is particularly vulnerable as it loops under the aortic arch), and the phrenic nerves, which run along the lateral aspects of the pericardium. Lymphatic drainage of the lungs and mediastinum is critical, with specific lymph node stations targeted during mediastinoscopy:

\begin{itemize}
  \item \textbf{Station 2R/L (Upper Paratracheal):} Located superior to the upper border of the aortic arch.
  \item \textbf{Station 4R/L (Lower Paratracheal):} Located inferior to the upper border of the aortic arch but superior to the lower border of the azygos vein (right) or the upper border of the pulmonary artery (left).
  \item \textbf{Station 7 (Subcarinal):} Located inferior to the carina, between the main bronchi.
\end{itemize}
Understanding the precise anatomical relationships of these nodes to the trachea, esophagus, and great vessels is paramount to safe and effective mediastinoscopy, minimizing the risk of vascular or neural injury.

\section*{Surgical Principle \& Rationale}
The fundamental surgical principle of mediastinoscopy involves creating a controlled, minimally invasive corridor into the pretracheal space through a small cervical incision. This allows for direct visualization and targeted biopsy of specific mediastinal lymph node stations and other mediastinal masses. A rigid mediastinoscope, equipped with a light source and magnification, is carefully advanced anterior to the trachea, providing a clear view of the anatomical landscape. The technical goal is to obtain representative and adequate tissue samples for definitive histological diagnosis while meticulously avoiding injury to vital surrounding structures.

The rationale for this procedure is rooted in the need for accurate tissue diagnosis and precise staging, particularly for non-small cell lung cancer (NSCLC). Imaging modalities like CT and PET scans can suggest mediastinal lymph node involvement (N2/N3 disease), but they lack the definitive diagnostic accuracy of histological examination. Mediastinoscopy provides this crucial histological confirmation, which is essential for guiding therapeutic decisions. For instance, confirmed N2/N3 disease often shifts treatment from primary surgical resection to neoadjuvant chemotherapy or definitive chemoradiation. Furthermore, for non-malignant conditions such as sarcoidosis or lymphoma, the biopsy provides the specific diagnosis necessary for targeted medical management, which cannot be achieved through imaging alone.

\section*{History \& Surgical Evolution}
The concept of mediastinal exploration dates back to the early 20th century, but it was Dr. Eric Carlens, a Swedish thoracic surgeon, who introduced cervical mediastinoscopy as a standardized diagnostic procedure in 1959. His innovative technique, involving a small suprasternal incision and the use of a rigid scope, revolutionized the staging of lung cancer and the diagnosis of other mediastinal diseases. Prior to Carlens' work, mediastinal assessment often required more invasive thoracotomy or blind needle biopsies, which carried higher morbidity and lower diagnostic yield.

Following Carlens' pioneering efforts, the technique rapidly gained acceptance worldwide, becoming the gold standard for surgical mediastinal staging of lung cancer throughout the latter half of the 20th century. Over time, advancements included the integration of video technology (video-assisted mediastinoscopy) to improve visualization, allow for team viewing, and facilitate documentation. Refinements in instrumentation, such as specialized biopsy forceps and improved light sources, further enhanced safety and efficacy. While newer, less invasive techniques like endobronchial ultrasound-guided transbronchial needle aspiration (EBUS-TBNA) have emerged and often serve as a first-line approach, mediastinoscopy continues to hold a crucial role, particularly when EBUS is non-diagnostic, technically challenging, or when a larger tissue sample is required for comprehensive pathological and molecular analysis.

\section*{Clinical Indications}
\begin{itemize}
  \item To accurately stage non-small cell lung cancer (NSCLC) by obtaining definitive histological samples from mediastinal lymph nodes, especially when imaging (CT, PET-CT) suggests N2/N3 disease.
  \item To diagnose suspected sarcoidosis by identifying non-caseating granulomas in mediastinal lymph nodes, particularly when other diagnostic avenues are inconclusive.
  \item To evaluate and diagnose suspected lymphoma or other primary mediastinal tumors, requiring a substantial tissue sample for subtyping and molecular analysis.
  \item To investigate unexplained mediastinal lymphadenopathy identified on imaging studies, where the etiology remains unclear after less invasive workup.
  \item To diagnose infectious processes such as tuberculosis or fungal infections affecting mediastinal nodes, which may require specific microbiological cultures.
  \item When less invasive diagnostic methods, such as endobronchial ultrasound-guided transbronchial needle aspiration (EBUS-TBNA), have been inconclusive, negative despite high clinical suspicion, or are technically not feasible.
  \item To assess resectability of lung tumors by confirming the absence of N2/N3 disease in patients considered for curative-intent surgical resection.
  \item To re-stage mediastinal lymph nodes after neoadjuvant therapy for lung cancer, though this role is more debated and often reserved for specific clinical scenarios.
\end{itemize}

\section*{Clinical Positioning}
Mediastinoscopy serves as both a primary and, in certain contexts, a secondary diagnostic and staging procedure. It is considered a primary surgical staging test for mediastinal disease, particularly in patients with non-small cell lung cancer where preoperative imaging (CT and PET scans) suggests involvement of mediastinal lymph nodes (N2/N3 disease). In such cases, mediastinoscopy provides the most definitive histological confirmation of nodal status, which is critical for guiding treatment decisions and determining surgical resectability. It is also a primary diagnostic tool for suspected mediastinal masses or lymphadenopathy of unknown etiology when a substantial tissue sample is required for diagnosis, especially when fine needle aspiration is insufficient.

However, with the advent and widespread availability of less invasive techniques like EBUS-TBNA, mediastinoscopy often functions as a secondary or confirmatory procedure. If EBUS-TBNA is negative for malignancy but there remains a high clinical suspicion of mediastinal nodal involvement, or if the samples obtained are insufficient for complete pathological analysis (e.g., for molecular testing or specific lymphoma subtyping), mediastinoscopy may be performed as a follow-up. It can also be considered secondary for re-staging after neoadjuvant therapy, though its role here is more limited and often debated due to the challenges posed by post-treatment fibrosis and inflammation.

\section*{Surgical Technique \& Procedure}
Mediastinoscopy is performed under general anesthesia, typically requiring endotracheal intubation to secure the airway. The patient is positioned supine with the neck hyperextended, often supported by a shoulder roll, to optimize exposure of the anterior neck and facilitate access to the superior mediastinum.

The procedure generally follows these steps:
\begin{enumerate}
  \item \textbf{Anesthesia and Positioning:} General anesthesia is induced, and the patient is intubated. The patient is placed in a supine position with the neck extended and a shoulder roll positioned to further hyperextend the neck, bringing the suprasternal notch forward.
  \item \textbf{Incision:} A small, approximately 3-4 cm transverse incision is made just above the suprasternal notch, typically within a skin crease for cosmetic reasons.
  \item \textbf{Dissection to Pretracheal Space:} The surgeon carefully dissects through the subcutaneous tissue and platysma muscle. The strap muscles (sternohyoid and sternothyroid) are identified and retracted laterally. The pretracheal fascia is then incised, and blunt dissection is performed to create a tunnel in the pretracheal space, anterior to the trachea and posterior to the innominate artery and aortic arch.
  \item \textbf{Mediastinoscope Insertion and Exploration:} Once the pretracheal space is adequately developed, the rigid mediastinoscope is inserted. The surgeon systematically explores the mediastinum, identifying key anatomical landmarks such as the anterior wall of the trachea, the innominate artery (which crosses anteriorly), the superior vena cava to the right, and the aortic arch to the left.
  \item \textbf{Lymph Node Identification and Biopsy:} Lymph nodes in stations 2R, 2L, 4R, 4L, and 7 are identified based on their anatomical relationships. These nodes are carefully dissected from surrounding tissues using specialized long, slender biopsy forceps. Biopsies are taken, ensuring adequate tissue for pathological analysis, and often sent for frozen section analysis to confirm diagnostic yield.
  \item \textbf{Hemostasis:} Meticulous hemostasis is achieved throughout the procedure using electrocautery, surgical clips, or topical hemostatic agents. Any bleeding must be controlled before proceeding or withdrawing the scope.
  \item \textbf{Scope Withdrawal and Wound Irrigation:} After satisfactory biopsies are obtained and hemostasis is confirmed, the mediastinoscope is slowly withdrawn. The surgical cavity is irrigated with saline to check for any residual bleeding.
  \item \textbf{Wound Closure:} The incision is then closed in layers. The strap muscles are reapproximated if necessary, the platysma is closed, and the skin is closed with absorbable sutures, non-absorbable sutures, or surgical staples.
  \item \textbf{Drains:} A small suction drain is rarely placed unless significant bleeding was encountered or a large cavity was created, which is uncommon for standard mediastinoscopy.
\end{enumerate}

\section*{Preoperative Workup \& Preparation}
Thorough preoperative preparation is essential to ensure patient safety and optimize the diagnostic yield of mediastinoscopy. This involves a comprehensive assessment of the patient's overall health and specific considerations related to the procedure.

\begin{itemize}
  \item \textbf{Medical Evaluation:} A comprehensive medical history and physical examination are performed, with particular attention to cardiopulmonary function, given the potential for underlying lung disease and the need for general anesthesia.
  \item \textbf{Laboratory Tests:} Routine blood tests are obtained, including a complete blood count (CBC), electrolyte panel, renal and liver function tests, and coagulation studies (PT/INR, PTT) to assess overall health and identify any bleeding risks.
  \item \textbf{Imaging Studies:} Recent computed tomography (CT) of the chest and often a positron emission tomography (PET-CT) scan are crucial. These images guide the surgeon to the most suspicious lymph node stations for biopsy, help identify potential anatomical challenges, and detect vascular anomalies.
  \item \textbf{Anesthesia Clearance:} Patients undergo a preoperative anesthesia evaluation to assess their fitness for general anesthesia, particularly if they have significant comorbidities such as severe cardiac or pulmonary disease.
  \item \textbf{Medication Review:} All medications are reviewed. Anticoagulants and antiplatelet agents (e.g., warfarin, aspirin, clopidogrel) typically need to be discontinued or bridged prior to surgery to minimize bleeding risk, following specific institutional protocols.
  \item \textbf{NPO Status:} Patients are instructed to be NPO (nil per os) for a specified period (usually 6-8 hours) before the procedure to prevent aspiration during anesthesia.
  \item \textbf{Informed Consent:} A detailed discussion of the procedure, its benefits, risks, alternatives, and potential need for further intervention (e.g., emergency sternotomy for bleeding) is conducted, and informed consent is obtained.
  \item \textbf{Antibiotic Prophylaxis:} Prophylactic antibiotics, typically a first or second-generation cephalosporin, are administered intravenously prior to the incision to reduce the risk of surgical site infection.
\end{itemize}

\section*{Contraindications \& Precautions}
\textbf{Absolute Contraindications:}
\begin{itemize}
  \item Uncorrectable coagulopathy or severe bleeding diathesis, which significantly increases the risk of catastrophic hemorrhage.
  \item Severe cervical spine immobility or fusion that prevents adequate neck hyperextension, making surgical access unsafe or impossible.
  \item Severe tracheal deviation or compression that would obscure anatomical landmarks or make scope insertion hazardous, increasing the risk of tracheal injury.
  \item Previous sternotomy with extensive scarring in the pretracheal space, which significantly increases the risk of major vascular injury due to obliterated tissue planes.
  \item Aneurysm of the aortic arch or great vessels that directly impinges on the surgical field, posing an extreme risk of rupture.
\end{itemize}

\textbf{Relative Contraindications \& Precautions:}
\begin{itemize}
  \item Prior mediastinoscopy, which can result in significant scarring and obliteration of tissue planes, making subsequent procedures more challenging and risky.
  \item Previous neck or mediastinal radiation therapy, which can also lead to fibrosis, altered anatomy, and increased tissue fragility.
  \item Superior vena cava (SVC) syndrome, as engorged and fragile veins can increase bleeding risk and obscure visualization.
  \item Major vessel encasement by tumor, particularly the innominate artery or aorta, which increases the risk of catastrophic hemorrhage during dissection or biopsy.
  \item Severe underlying cardiopulmonary disease, which may increase anesthetic risks and compromise the patient's ability to tolerate the procedure.
  \item Significant atherosclerotic disease of the great vessels, requiring extreme caution during dissection to avoid plaque dislodgement and stroke.
  \item Significant obesity or short, thick neck, which can make anatomical identification and dissection more difficult and prolong operative time.
\end{itemize}

\section*{Complications \& Risks}
While generally safe, mediastinoscopy carries potential risks, some of which can be serious and life-threatening. These can be broadly categorized into general surgical risks and procedure-specific complications.

\begin{itemize}
  \item \textbf{General Surgical Complications:}
    \begin{itemize}
      \item \textbf{Bleeding:} Minor bleeding from small vessels is common and usually controlled with electrocautery. Major hemorrhage can occur from injury to large mediastinal vessels (e.g., innominate artery, aorta, pulmonary artery, superior vena cava), potentially requiring emergency sternotomy or thoracotomy.
      \item \textbf{Infection:} Surgical site infection (wound infection) at the incision site. Mediastinitis, a rare but severe infection of the mediastinal space, can be life-threatening.
      \item \textbf{Anesthesia Complications:} Adverse reactions to anesthetic agents, cardiac events (e.g., arrhythmia, myocardial infarction), respiratory complications (e.g., aspiration, pneumonia), or prolonged intubation.
    \end{itemize}
  \item \textbf{Procedure-Specific Risks:}
    \begin{itemize}
      \item \textbf{Recurrent Laryngeal Nerve Injury:} Can cause temporary or permanent hoarseness or vocal cord paralysis. The left recurrent laryngeal nerve is particularly vulnerable due to its longer course and proximity to the aortic arch.
      \item \textbf{Pneumothorax:} Collapse of a lung due to injury to the pleura, which may require chest tube insertion for re-expansion.
      \item \textbf{Esophageal Injury:} Rare but serious, can lead to mediastinitis, fistula formation, or stricture.
      \item \textbf{Tracheal Injury:} Rare, but can lead to air leak, mediastinitis, or tracheoesophageal fistula.
      \item \textbf{Chylothorax:} Extremely rare, due to injury to the thoracic duct, leading to lymphatic fluid accumulation in the chest.
      \item \textbf{Air Embolism:} Rare, but potentially fatal, if air enters a major mediastinal vein during dissection.
      \item \textbf{Stroke:} Very rare, due to manipulation near the great vessels, dislodgement of atherosclerotic plaque, or pre-existing vascular disease.
      \item \textbf{Injury to the Phrenic Nerve:} Can cause diaphragmatic paralysis, leading to respiratory compromise, though very rare.
      \item \textbf{Inadequate Biopsy:} Failure to obtain sufficient or representative tissue for diagnosis, necessitating further diagnostic procedures.
      \item \textbf{Mediastinal Hematoma:} Accumulation of blood in the mediastinal space, potentially causing airway compression.
    \end{itemize}
\end{itemize}

\section*{Postoperative Recovery Timeline}
Immediately after mediastinoscopy, the patient is transferred to the Post-Anesthesia Care Unit (PACU) for close monitoring as they recover from general anesthesia. Vital signs, oxygen saturation, and pain levels are continuously assessed. The surgical site is checked for bleeding or swelling, and the patient's voice is often evaluated for any hoarseness, which could indicate recurrent laryngeal nerve irritation or injury. Most patients experience mild to moderate neck pain and soreness, which is effectively managed with oral analgesics. Early ambulation is encouraged to prevent complications such as deep vein thrombosis.

Within the first 24-48 hours, most patients are stable enough for discharge, often staying overnight for observation. They are encouraged to resume a light diet as tolerated and gradually increase their activity. Instructions are provided for wound care, pain management, and activity restrictions, typically advising against heavy lifting or strenuous activities for one to two weeks. Patients are monitored for signs of complications such as increasing pain, fever, difficulty breathing, persistent hoarseness, or redness/swelling at the incision site. Long-term recovery usually involves complete resolution of neck discomfort and healing of the incision site, leaving a small, usually inconspicuous scar. Full return to normal activities is typically expected within 2-3 weeks.

\section*{Surgical Findings \& Pathology}
During mediastinoscopy, the surgeon documents visual findings such as the size, consistency (e.g., firm, rubbery), and appearance (e.g., inflammatory, necrotic) of the lymph nodes or masses encountered. Enlarged, firm, or matted lymph nodes are often suspicious for malignancy or granulomatous disease. The extent of nodal involvement and any adherence to surrounding structures are also noted. These intraoperative findings guide the selection of biopsy sites and provide context for the subsequent pathological analysis.

The pathology report, derived from the biopsied tissue, provides the definitive diagnosis and is the most critical outcome of the procedure. For suspected lung cancer, the pathologist will determine the presence or absence of malignant cells, the specific histological subtype of non-small cell lung cancer (e.g., adenocarcinoma, squamous cell carcinoma), and whether the cancer has spread to the biopsied lymph nodes. This determines the N-status (nodal involvement) for staging. For non-cancerous conditions, the report might reveal specific findings such as non-caseating granulomas characteristic of sarcoidosis, specific cellular patterns indicative of lymphoma, or evidence of infectious processes like tuberculosis. The adequacy of the tissue sample for diagnosis and any additional molecular or immunohistochemical testing is also reported.

\section*{Expected Postoperative Course}
A normal, successful postoperative course following mediastinoscopy typically involves a smooth recovery from general anesthesia with minimal discomfort. Patients can expect mild to moderate neck pain and soreness at the incision site, which is usually well-controlled with oral pain medications. Hoarseness, if present, is often transient due to irritation of the recurrent laryngeal nerve during dissection and usually resolves within days to a few weeks. Patients are typically able to resume a light diet within hours of the procedure and are encouraged to ambulate early.

The incision site is expected to heal cleanly, with minimal redness or swelling, and the sutures or staples are usually removed at a follow-up visit within 1-2 weeks. Most patients can return to light activities within a few days and to their full normal routine, including work and exercise, within 2-3 weeks, provided there are no complications. The absence of fever, significant bleeding, persistent hoarseness, or difficulty breathing indicates an uneventful recovery. The primary focus shifts to awaiting the definitive pathology results, which will guide any subsequent treatment decisions.

\section*{Postoperative Care \& Follow-up}
Postoperative care and follow-up are crucial for monitoring recovery, discussing results, and planning subsequent treatment.

\begin{itemize}
  \item \textbf{Postoperative Clinic Visit:} Typically scheduled 1-2 weeks after the procedure to check the incision site, remove sutures or staples if present, and discuss the final pathology report in detail.
  \item \textbf{Pathology Review:} The surgeon will explain the implications of the biopsy results for diagnosis, staging, and treatment options, ensuring the patient understands the next steps.
  \item \textbf{Multidisciplinary Team Discussion:} For cancer patients, the results are often discussed in a multidisciplinary tumor board involving oncologists, radiation oncologists, and thoracic surgeons to formulate the optimal, individualized treatment plan.
  \item \textbf{Further Imaging/Labs:} Depending on the pathology findings and the subsequent treatment plan, further imaging (e.g., repeat CT, MRI) or laboratory tests may be ordered.
  \item \textbf{Activity Resumption:} Patients are advised on a gradual return to normal activities, typically avoiding strenuous exercise or heavy lifting for several weeks to allow for complete wound healing.
  \item \textbf{Warning Signs:} Patients are instructed to contact their doctor immediately if they experience fever (above 101$^{\circ}$F or 38.3$^{\circ}$C), worsening neck pain, significant redness or swelling at the incision site, pus-like drainage, persistent or worsening hoarseness, difficulty breathing, or any other concerning symptoms.
\end{itemize}
Adherence to these follow-up instructions is vital for ensuring a complete recovery and effective management of the underlying condition.

\section*{Epidemiology}
Mediastinoscopy is predominantly performed in the context of lung cancer, which remains a leading cause of cancer-related mortality worldwide. Non-small cell lung cancer (NSCLC) accounts for approximately 85\% of all lung cancers, and accurate mediastinal lymph node staging is paramount for treatment selection and prognosis. The incidence of lung cancer is higher in older adults, typically peaking in the 60s and 70s, and is slightly more common in men than women, although rates in women are increasing. The procedure's frequency is directly tied to the incidence of lung cancer and the prevalence of other mediastinal pathologies like sarcoidosis and lymphoma. While EBUS-TBNA has reduced the overall number of mediastinoscopies performed, it remains a vital tool, especially in regions where EBUS is unavailable, when EBUS results are inconclusive, or when larger tissue samples are required. The morbidity rate for mediastinoscopy is generally low, ranging from 1-5\%, with major complications like significant hemorrhage or recurrent laryngeal nerve injury occurring in less than 1-2\% of cases. Mortality is exceedingly rare, typically less than 0.1-0.2\%, primarily associated with catastrophic vascular injury.

\section*{Outcomes \& Prognosis}
The primary outcome of mediastinoscopy is the provision of a definitive histological diagnosis and accurate mediastinal lymph node staging. For lung cancer, the diagnostic accuracy for mediastinal nodal involvement (N2/N3 disease) is very high, often exceeding 90-95\%, making it a highly reliable staging tool. This accuracy is crucial for determining the resectability of lung tumors and guiding the choice between surgery, neoadjuvant therapy, or definitive chemoradiation. Patients found to have N2 or N3 disease typically have a poorer prognosis compared to those with N0 or N1 disease, and their treatment pathways are significantly altered. For example, the APPAC trial demonstrated the efficacy of neoadjuvant chemotherapy in locally advanced NSCLC, where accurate staging by mediastinoscopy is critical.

For non-malignant conditions like sarcoidosis or lymphoma, mediastinoscopy provides the tissue diagnosis necessary to initiate appropriate medical therapy, leading to improved functional outcomes and disease control. The overall prognosis following mediastinoscopy is largely dictated by the underlying disease process rather than the procedure itself. Successful staging or diagnosis allows for tailored treatment, which directly impacts patient survival and quality of life. Recurrence of mediastinal disease, if applicable, is a function of the primary pathology and its response to treatment, rather than a direct outcome of the mediastinoscopy. Prognostic factors for lung cancer patients undergoing mediastinoscopy include the number and location of positive lymph nodes, tumor histology, and overall patient health status.

\section*{References}
\begin{enumerate}
  \item MedlinePlus. Mediastinoscopy. U.S. National Library of Medicine; 2024. Available from: \url{https://medlineplus.gov/ency/article/003861.htm}
  \item Sabiston DC, Townsend CM. Sabiston Textbook of Surgery: The Biological Basis of Modern Surgical Practice. 21st ed. Elsevier; 2021.
  \item National Comprehensive Cancer Network (NCCN). NCCN Clinical Practice Guidelines in Oncology: Non-Small Cell Lung Cancer. Version 5.2024. Available from: \url{https://www.nccn.org/professionals/physician_gls/pdf/nscl.pdf}
  \item Carlens E. Mediastinoscopy: a new method for inspection and tissue biopsy in the superior mediastinum. Dis Chest. 1959;36(3):343-352.
  \item De Leyn P, et al. ESTS guidelines for intraoperative lymph node staging in non-small cell lung cancer. Eur J Cardiothorac Surg. 2014;45(5):787-801.
\end{enumerate}

\clearpage